\documentclass[../../../main/main.tex]{subfiles}

\begin{document}

\chapter{Axioms}

\section{Axioms of Equality}
\begin{axiom}[Idenity/Reflexivity of Equality]
    \[
    \forall A \in \mathbf{P}, A = A.
    \]
\end{axiom}
\begin{axiom}[Symmetry of Equality]
    \[
    \forall A, B \in \mathbf{P}, A = B \iff B = A.
    \]
\end{axiom}
\begin{axiom}[Transitivity of Equality]
    \[
    \forall A, B, C \in \mathbf{P}, A = B \land B = C \implies A = C.
    \]
\end{axiom}

\section{Axioms of Betwneenness}
The betweeness relation is governed by the axioms of betweeness.
\begin{axiom}[Idenity/Reflexivity of Betweenness]
    \[
    \forall A, B \in \mathbf{P}, \mathcal{B}(A, B, A) \iff A = B.
    \]
\end{axiom}
\begin{axiom}[Symmetry of Betweenness]
    \[
    \forall A, B, C \in \mathbf{P}, \mathcal{B}(A, B, C) \iff \mathcal{B}(C, B, A).
    \]
\end{axiom}
\begin{axiom}[Order of Betweeness]
    \[
    \forall A, B, X \in \mathbf{P}, \mathcal{B}(A, B, X) \implies (\neg \mathcal{B}(B, A, X) \land \neg \mathcal{B}(A, X, B)).
    \]
\end{axiom}
\begin{axiom}[Inner Point/Density of Betweeness]
    \[
    \forall A, C \in \mathbf{P}, A \neq C \implies \exists B \in \mathbf{P} \colon \mathcal{B}(A, B, C).
    \]
\end{axiom}
\begin{axiom}[Outer Point/Extensibility of Betweenness]
    \[
    \forall A, B \in \mathbf{P}, A \neq B \implies \exists C \in \mathbf{P} \colon \mathcal{B}(A, B, C).
    \]
\end{axiom}

\section{Axioms of Congruence}
\[ \mathcal{S} \coloneqq \{ S \subseteq \mathbf{P} \mid 1 \le |S| \le 2 \}. \]
\[ \cong \subseteq \mathcal{S} \times \mathcal{S}. \]
\begin{axiom}[Reflexivity of Congruence]
    \[
    \forall S \in \mathcal{S}, S \cong S.
    \]
\end{axiom}
\begin{axiom}[Symmetry of Congruence]
    \[
    \forall S_1, S_2 \in \mathcal{S}, S_1 \cong S_2 \iff S_2 \cong S_1.
    \]
\end{axiom}
\begin{axiom}[Transitivity of Congruence]
    \[
    \forall S_1, S_2, S_3 \in \mathcal{S}, (S_1 \cong S_2 \land S_2 \cong S_3) \implies S_1 \cong S_3.
    \]
\end{axiom}
\begin{axiom}[Idenity of Indiscernibles]
    \[
    \forall S_1, S_2 \in \mathcal{S}, (|S_1| = 1 \land S_1 \cong S_2 ) \implies |S_2| = 1.
    \]
\end{axiom}

\section{idk}
The ruler axiom states that given a point-pair $S$ and two discint points $A$ and $B$, there exists one point $X$ such that $\mathcal{B}(A, B, X)$ and $\{B, X\}$ is congruent to $S$.
\begin{axiom}[Ruler Axiom]
    \[
    \forall S \in \mathcal{S}, \forall A, B \in \mathbf{P}, A \neq B \implies \exists!X \in \mathbf{P} \colon \mathcal{B}(A, B, X) \land \{B, X\} \cong S.
    \]
\end{axiom}

\end{document}